\section{Prática com Git}

\begin{frame}{Nosso projeto}

  \begin{block}{Objetivo}
    Vamos escrever o cardápio de um restaurante utilizando o Git para acompanhar o processo de desenvolvimento.
  \end{block}
  \begin{columns}[T]
    \begin{column}{0.48\textwidth}
      \begin{block}{Ambiente Local}
        No seu computador:
        \begin{itemize}
          \item Instale o Git pelo site oficial.
          \item Crie uma conta no GitHub.
        \end{itemize}
        \vspace{0.2cm}
        \centering

        \href{https://git-scm.com/}{\textbf{🌐 git-scm.com}} \quad \href{https://github.com/}{\textbf{🌐 github.com}}
      \end{block}
    \end{column}
    \begin{column}{0.48\textwidth}
      \begin{block}{Ambiente Online}
        Vamos utilizar o GitHub Codespaces.

        \vspace{0.2cm}
        Ele fornece um ambiente de desenvolvimento já configurado para trabalharmos com o Git.

        \vspace{0.2cm}
        \centering
        \href{https://github.com/codespaces}{\textbf{🌐 github.com/codespaces}}
      \end{block}
    \end{column}
  \end{columns}
\end{frame}

\begin{frame}{Primeiro passo: identificar o autor}
  \begin{block}{Configurando o Git}
    Antes de começar, vamos informar quem está fazendo os commits.

    \vspace{0.3cm}

    \texttt{git config --global user.name "Your Name"}

    \texttt{git config --global user.email "you@example.com"}
  \end{block}
  No \textbf{GitHub Codespaces}, essa configuração já vem preparada para você.
\end{frame}


\begin{frame}{Criando o projeto}

  \begin{block}{Nosso projeto}
    Vamos criar o repositório do cardápio da
    \textbf{Cantina da Poli}.
  \end{block}

  \begin{block}{Criando o diretório}
    \texttt{mkdir cantina-da-poli}

    \texttt{cd cantina-da-poli}
  \end{block}

  \begin{block}{\texttt{git init} — inicia o repositório}
    \texttt{git init}
  \end{block}

\end{frame}


\begin{frame}{Nosso primeiro commit}

  \begin{block}{Criando o cardápio}
    Vamos criar o primeiro arquivo do projeto:

    \vspace{0.3cm}
    \texttt{index.html}
  \end{block}

  \begin{block}{Verificando as alterações}
    \texttt{git status}
  \end{block}

  \begin{block}{Preparando o commit}
    \texttt{git add .}
  \end{block}

  \begin{block}{Registrando o primeiro estado do projeto}
    \texttt{git commit -m "Cria o cardápio da Cantina da Poli"}
  \end{block}

\end{frame}


\begin{frame}{A Cantina está crescendo}

  \begin{block}{Adicionando uma nova funcionalidade}
    Vamos alterar o cardápio e adicionar novos pratos.
  \end{block}

  \begin{block}{Verificando o que mudou}
    \texttt{git status}
  \end{block}

  \begin{block}{Registrando a alteração}
    \texttt{git add .}

    \texttt{git commit -m "Adiciona novos pratos ao cardápio"}
  \end{block}

\end{frame}


\begin{frame}{Consultando a história do projeto}

  \begin{block}{\texttt{git log} — histórico de commits}
    \texttt{git log}

    \vspace{0.2cm}
    Podemos ver como o cardápio foi evoluindo ao longo do tempo.
  \end{block}

  \begin{block}{\texttt{git show} — detalhes de um commit}
    \texttt{git show}

    \vspace{0.2cm}
    Mostra as alterações realizadas em um commit.
  \end{block}

\end{frame}


\begin{frame}{Publicando o projeto}

  \begin{block}{Repositório remoto}
    Agora queremos compartilhar o projeto da
    \textbf{Cantina da Poli} com outras pessoas.
  \end{block}

  \begin{block}{\texttt{git push} — envia os commits}
    \texttt{git push}
  \end{block}

\end{frame}

\begin{frame}{Arquivos úteis}
  \begin{itemize}
    \item \texttt{.gitignore} — listar arquivos que não devem ser versionados
    \item \texttt{.gitattributes} — controlar tratamento de arquivos, diffs, merges
  \end{itemize}
\end{frame}
